\documentclass[a4paper,12pt]{article}
\usepackage{color}
\usepackage{graphicx, gensymb}
\usepackage{amsfonts, cancel, amssymb}
\usepackage{amsmath, array, esvect, esint}
\usepackage{typearea}
\usepackage[a4paper,left=2cm,right=2cm,top=2cm,bottom=3cm,bindingoffset=5mm]{geometry}

\begin{document}

\title{Leiterschleifen und Induktion}
\author{Joachim Schwardt}
\date{\today}
\maketitle

\section{$A=A(t)$}
\subsection*{a)}
Induktionsgesetz: 
\begin{align*}
U_{ind}=-\frac{d}{dt}\iint_A \vv{B}\cdot d\vv{A} \overset{\substack{ \vv{B}=const\\|}}{=}-B\frac{d}{dt}\iint_A \cos(\varphi(t))dA = a^2B\omega\sin(\omega t)
\end{align*}
\subsection*{b)} Kraft auf Ströme:
\begin{align*}
d\vv{F}&=q\vv{v}\times\vv{B}=Idt\vv{v}\times\vv{B}=I\vv{dl}\times\vv{B}\\
&\Rightarrow\ F = IlB\sin(\varphi)
\end{align*}
Die mechanische Leistung entspricht dann:
\begin{align*}
P = vF = \omega a\ q|\vv{v}\times\vv{B}|=\omega a\ IaB\sin(\omega t)=\frac{a^4B^2\omega^2\sin^2(wt)}{R}
\end{align*} 
\section{$B = B(t)$}
\subsection*{a)}
Nach Maxwell berechnet sich die Flussdichte des Leiters folgendermaßen:
\begin{align*}
\oint_{\partial A} \vv{B}\cdot\vv{ds} = \mu_0I(t) = \mu_0I_0\sin(\omega t)\ \Rightarrow\ B(t) = \frac{\mu_0I_0\sin(\omega t)}{2\pi r}
\end{align*}
Da $\vv{B}\bot\vv{A}$ folgt aus dem Induktionsgesetz:
\begin{align*}
U_{ind}&=-\dot{\Phi}_m = -\frac{d}{dt}\iint_A B(t)dA = -\frac{d}{dt}\int_0^{a}\int_b^{b+a} \frac{\mu_0I_0\sin(\omega t)}{2\pi x} dxdy=\\
&=-\frac{d}{dt}\frac{a\mu_0I_0\sin(\omega t)}{2\pi}\mathrm{ln}\left(1+\frac{a}{b}\right) = -\frac{a\omega\mu_0I_0\cos(\omega t)}{2\pi}\mathrm{ln}\left(1+\frac{a}{b}\right)
\end{align*}
\subsection*{b)}
Die Kraft auf die Leiterschleife beträgt:
\begin{align*}
\vv{F} =I_{ind}\vv{l}\times\vv{B}(t) = \frac{I_{ind}\mu_0I_0\sin(\omega t)}{2\pi}\left(\frac{a}{b}-\frac{a}{b+a}\right)\vv{e_x}=\frac{\mu_0}{2\pi}I_{ind}I_0\sin(\omega t)\left(\frac{a^2}{b(a+b)}\right)\vv{e_x}
\end{align*}
\subsection*{c)}
Aus a) kennen wir bereits $\Phi_m$. Damit ergibt sich:
\begin{align*}
\Phi_m = L_{12}I(t)\ \Rightarrow\ L_{12} = \frac{a\mu_0I(t)}{2\pi I(t)}\mathrm{ln}\left(1+\frac{a}{b}\right) = \frac{a\mu_0}{2\pi}\mathrm{ln}\left(1+\frac{a}{b}\right)
\end{align*}
\section{Beweglicher Leiter im homogenen Magnetfeld} 
\subsection*{a)}
Die Induktionsspannung aufgrund der Flächenänderung muss berücksichtigt werden:
\begin{align*}
U=U_0+U_{ind}=U_0-\frac{d}{dt}\iint_A \vv{B}\cdot d\vv{A}=U_0-B\frac{d}{dt}xd=U_0-B\dot{x}d
\end{align*}
Newton'sche Bewegungsgleichung und Kraft auf einen stromdurchflossenen Leiter:
\begin{align*}
F = m\ddot{x}\ \Rightarrow\ \dot{v}=\frac1m\ IBd = \frac1m\ \frac{U}{R}Bd = \frac{BdU_0}{Rm}-\frac{B^2d^2v}{Rm} 
\end{align*}
\subsection*{b)}
Die Lösung dieser Differentialgleichung ergibt:
\begin{align*}
v(t)&=v_h(t)+v_p(t)=C_1e^{-\frac{B^2d^2t}{Rm}}+\frac{U_0}{Bd}\overset{\substack{v(t=0)\overset{!}{=}0 \\|}}{=}\frac{U_0}{Bd}\left(1-e^{-\frac{B^2d^2t}{Rm}}\right)\\
&\Rightarrow\ v_{max} = \frac{U_0}{Bd}
\end{align*}
Alternativ kann einfach die Formel für Spannung aus 3a) betrachtet werden:
\begin{align*}
U = U_0-Bdv\overset{!}{=}0\ \Rightarrow\ v_{max} = \frac{U_0}{Bd}
\end{align*}
\section{Kreisförmige Leiterschleife in $B=B(t)$}
\begin{align*}
U_{ind}(t)&=-N\dot{\Phi}_m(t)=-N\frac{d}{dt}B(t)\cos(\alpha)\pi R^2=N\frac{B_0}{\tau}\pi R^2\cos(\alpha)e^{-\frac{t}{\tau}}\\
&\Rightarrow\ U_{ind}(t_1=2s) = -10\frac{2T}{1s}\frac{\pi}{e^2} (0.04m)^2\cos(45\degree) \simeq 9.62mV
\end{align*}
\end{document}